\BeleskeID{03-s020b}
% element: 03-s020b:e03
% element: 03-s020b:e04
% element: 03-s020b:d02

Neka je \(H=I_3+I_2+I_1+I_0\) lokalno prisustvo zahteva. Pošto je u prikazanoj realizaciji \(GS=EI\,H\), izraz za prenos dozvole može se napisati i kao
\[
 EO=EI\,\overline{GS}=EI\,\overline H.
\]
Poslednja jednakost sledi razmatranjem EI: kada je nula, EO je nula nezavisno od H; kada je jedinica, GS je jednak H. Kada rad nije dozvoljen, oba stanja lokalnog zahteva daju isti EO, jednak nuli.

Ako je \(EI=1,H=0\), koder sme da radi, ali nema šta da koduje; \(EO=1\) dozvoljava rad sledećem. Ako je \(EI=1,H=1\), obrađuje lokalni zahtev i sa \(EO=0\) sprečava niže kodere da doprinesu rezultatu. Ako je \(EI=0\), već je izgubio dozvolu zbog višeg prioriteta i mora preneti tu zabranu dalje, čak i kada sam nema zahtev.

Na primer, neka viša grupa ima zahtev, srednja nema, a niža ima. Srednja dobija \(EI=0\). Inverzija njenog uslovljenog GS bila bi jedan i pogrešno bi dozvolila nižu grupu. Faktor EI u EO zadržava zabranu i kroz tu praznu srednju grupu.

Moguće je i da GS ne bude uslovljen EI, nego da označava samo H. To ne smeta izboru grupe na drugom nivou kodera prioriteta, koji sam bira najveći indeks među aktivnim grupama. Međutim, ni takav \(\overline{GS}\) nije dovoljan za EO: grupa bez lokalnog zahteva mogla je već izgubiti dozvolu zbog zahteva u višoj grupi. I tada je za prenos dozvole potreban uslov \(EI\,\overline H\).
